\documentclass[12pt,a4paper]{article}
\usepackage[UTF8]{ctex}
\usepackage{amsfonts,amsmath,amsthm,amssymb}
\usepackage{sfmath}
\usepackage{hyperref}
\renewcommand{\familydefault}{\sfdefault}
\usepackage{geometry}
\geometry{left=0cm,right=0cm,top=0cm,bottom=0.1cm, footskip=0cm, includefoot=true}
%\usepackage[contents = \thepage, color=yellow, angle=0, scale = 20]{background}
\pagestyle{empty}
%\usepackage{fancyhdr}
%\pagestyle{fancy}
%\fancyhead{}
%\fancyfoot{}
%\cfoot{}
%\lfoot{{\small \color{red} \slshape \thepage 页}}
%\fancyfoot[R]{\thepage}

\usepackage[all]{xy}
\usepackage{color}

\renewenvironment{proof}{{\noindent\it\textbf{Proof}}\quad\it}{\hfill$\square$}

%\theoremstyle{definition}

\newtheorem{thm}[equation]{Theorem}
\newtheorem{cor}[equation]{Corollary}
\newtheorem{lem}[equation]{Lemma}
\newtheorem{prop}[equation]{Proposition}
\newtheorem{defn}[equation]{Definition}
\newtheorem{rem}[equation]{Remark}
\newtheorem{conj}[equation]{Conjecture}

\newcommand{\Z}{\mathbb{Z}}

\title{}
\author{}


\begin{document}

\sffamily
\maketitle
\thispagestyle{empty}

\abstract{}
\tableofcontents

\large

\begin{sloppy}
\clubpenalty=-100
\widowpenalty=-100
 \section{extension SS} \label{ess}
Let $f$ be a morphism of converging spectral sequences
$$f_E: V_1\rightarrow V_2$$
这是谱序列的态射
$$f_A:A_1\rightarrow A_2$$
这是收敛到的群的态射

\paragraph{(1)} 将$f_A$看作只有两项的filtered complex

\paragraph{(2)} Extension谱序列定义为上述filtered complex诱导的谱序列

\paragraph{(3)} 证明Extension谱序列的$E_0$页同构于V的$E_\infty$页

\paragraph{(4)}证明如下结论：


\begin{defn}[Notation 2.4]
For $x \in E_\infty^{s}(V_1)$（这个notation表示在filtration degree s）, we denote by $\{x\}$ the set of all classes in $A_1$ that are detected by $x$. We use $[x]$ to refer to a specific class or a general class in $\{x\}$, with the choice understood from the context.
\end{defn}

\begin{prop}[Proposition 2.5]
Consider $f: V_1 \to V_2$, $x \in E_\infty^{s}(V_1)$, $y \in E_\infty^{s+n}(V_2)$ and $y' \in E_\infty^{s+m}(V_2)$ for $m, n \geq 0$.
\begin{enumerate}
\item[(1)] There is an $f$-extension from $x$ to $y$, i.e., $d_n^f(x) = y$ in the $f$-ESS if and only if there is a class $[x] \in \{x\}$ such that $f[x] \in \{y\}$.
\item[(2)] An $f$-extension from $x$ to $y$ is inessential, i.e., $y$ is trivial in the $E_n$-page of the $f$-ESS if and only if there exists an element $x' \in E_\infty^{s+a}(V_1)$ for $0 < a \leq n$ such that we have an essential differential $d_{n-a}^f(x') = y$. Equivalently, there exists a class $[x'] \in \{x'\} \subset A_1$ with $\mathrm{Fil}(x') > \mathrm{Fil}(x)$, such that $f[x'] \in \{y\}$.
\item[(3)] Suppose we have an $f$-extension from $x$ to both $y$ and $y'$.
\begin{enumerate}
\item[(a)] If $m = n$, then $y - y' \in {{}^f\!B}_{n-1}^{s+n}(V_2)$ and there exists an element $x' \in E_\infty^{s+a}(V_1)$ for $0 < a \leq n$ such that we have an essential differential $d_{n-a}^f(x') = y - y'$ in $f$-ESS. Equivalently, there exists a class $[x'] \in \{x'\} \subset A_1$ with $\mathrm{Fil}(x') > \mathrm{Fil}(x)$, such that $f[x'] \in \{y - y'\}$.
\item[(b)] If $m > n$, then the $f$-extension from $x$ to $y$ is inessential.
\end{enumerate}
\end{enumerate}
\end{prop}
\begin{proof}
Part (1) is straightforward from the setup of the extension spectral sequence and Definition。 Parts (2) and (3) follow from Part (1).
\end{proof}

\paragraph{(5)} 根据一般的理论建立estension谱序列中微分crossing的相关定义

\paragraph{(6)}证明：

\begin{rem}[Remark 2.8]
The following statements are equivalent.
\begin{itemize}
\item An $f$-extension $d_n^f(x) = y$ has no crossing that hits the range $\mathrm{Fil} \geq p$.
\item If for any $a > 0$ there is an $f$-extension from $x' \in E_\infty^{s+a}(V_1)$ to some nontrivial $y'$, then $\mathrm{Fil}(y') < p$ or $\mathrm{Fil}(y') > \mathrm{Fil}(y)$.
\end{itemize}
\end{rem}

\paragraph{(7)}证明：

\begin{prop}[Proposition 2.10]
An $f$-extension from $x$ to $y$ has no crossing in the range $\mathrm{Fil} \geq p$ if and only if for all $[x] \in \{x\}$ such that $\mathrm{Fil}(f[x]) \geq p$ we have $f[x] \in \{y\}$. In particular, we have
\begin{enumerate}
\item[(1)] An $f$-extension from $x$ to $y$ has no crossing if and only if for all $[x] \in \{x\}$ we have $f[x] \in \{y\}$.
\item[(2)] An $f$-extension $d_n^f(x) = 0$ has no crossing if and only if for all $[x] \in \{x\}$,
$$\mathrm{Fil}(f[x]) > \mathrm{Fil}(x) + n.$$
\end{enumerate}
\end{prop}
\begin{proof}
First we prove the only if part. Assume by contradiction that there exists $[x] \in \{x\}$ such that $f[x] = [y'] \notin \{y\}$ is detected by $y' \in E_\infty(V_2)$ with $\mathrm{Fil}(y') \geq p$. Then there is an $f$-extension from $x$ to $y'$. By Proposition 2.5 (3), $y'$ (if $\mathrm{Fil}(y) > \mathrm{Fil}(y')$) or $y - y'$ (if $\mathrm{Fil}(y) = \mathrm{Fil}(y')$) should be hit by a shorter $d^f$-differential. This shorter differential is a crossing of the $f$-extension from $x$ to $y$ that hits $\mathrm{Fil} \geq p$, which is a contradiction. Now we prove the if part. Assume that the $f$-extension has a crossing from $x'$ to $y'$ with $p \leq \mathrm{Fil}(y') \leq \mathrm{Fil}(y)$. Then there exists $[x'] \in \{x'\}$ such that $f([x']) \in \{y'\}$. Note that $[x] + [x'] \in \{x\}$ since $\mathrm{Fil}(x') > \mathrm{Fil}(x)$, we have
$$f([x] + [x']) = f([x]) + f([x']) \in \begin{cases} \{y'\} & \text{if } \mathrm{Fil}(y') < \mathrm{Fil}(y), \\ \{y + y'\} & \text{if } \mathrm{Fil}(y') = \mathrm{Fil}(y). \end{cases}$$
However, by assumption we should have $f([x] + [x']) \in \{y\}$ and it is contradictory to the equation above in either case.

Hence the main statement is true. Part (1) (2) are special cases of the main statement and Part (3) holds for degree reasons.
\end{proof}

\paragraph{(8)}证明如下结论，可以根据我们已有的谱序列收敛性的条件加相关的有界性或mittag-leffler条件：

\begin{prop}[Proposition 2.20]
Suppose that the sequence of converging spectral sequences
$$V_1 \xrightarrow{f} V_2 \xrightarrow{g} V_3$$
induces an exact sequence of target groups
$$A_1 \xrightarrow{\pi_*f} A_2 \xrightarrow{\pi_*g} A_3$$
that is exact in $A_2$. Then all permanent cycles in $E_\infty^{*,*}(V_2)$ of the $g$-ESS are boundaries in the $f$-ESS.
\end{prop}
\begin{proof}
If we consider the following sequence
$$0 \to \pi_*V_1 \xrightarrow{\pi_*f} \pi_*V_2 \xrightarrow{\pi_*g} \pi_*V_3 \to 0$$
and treat it as a chain complex filtered by the Adams filtration, we obtain an extension spectral sequence comprising $d^f$ and $d^g$ differentials. The abutment of this spectral sequence, when projected onto the $V_2$ component, is zero. Therefore, all permanent $d^g$-cycles are killed by $d^f$-differentials.

Again we offer another proof via homotopy groups. Assume that $y \in E_\infty^{*,*}(V_2)$ is a permanent $d^g$-cycle. By the convergence of ESS we know that it detects a $g$-torsion $[y] \in \{y\}$ with $g([y]) = 0$. Hence there exists $[x] \in \pi_*V_1$ such that $f([x]) = [y]$ by exactness. Assume that $[x]$ is detected by $x \in E_\infty^{*,*}(V_1)$. Then there exists an $f$-extension from $x$ to $y$.
\end{proof}

\section{修改意见}

\paragraph{} ExtensionSS 里面，\begin{enumerate}
    \item 删除 assDiff和essBoundary
    \item 增加 underlyingComplex,定义为，$A_1\rightarrow A_2$，viewed as filtered chain complex 
    \item 增加 ess，定义为 FilteredComplex.toSpectralSequence underlyingComplex, 可以加bounded条件
\end{enumerate}

\paragraph{}FourSpectraChain的名字改成ThreeSpectraChian，并且
\begin{enumerate}
    \item 增加 $cm_{13}$ 定义为 $cm_12$ composed with $cm_{23}$
    \item 增加 $essH = ExtensionSS\; cm_{13}$
\end{enumerate}

\paragraph{}
￥$ess\_exists$ 调用FilteredComplex的定理，可以假设bounded

\paragraph{}
$ess\_boundary\_zero$彻底删除，改为：
\paragraph{} $ess\_boundary\_nesteing$ 和 $ess\_boundary_subobject$ 删除

\paragraph{} $ess\_e0\_iso$删除，改为：
\begin{enumerate}
    \item 证明$ess$的$E_0$page 标准同构于 $E_1$的$E_\infty$和$E_2$的$E_\infty$的直和
    \item 证明ess的微分只有$E_1.page \infty$(的sub)到$E_2.page \infty$（的quotient）的分量，其他三个分量为0
    \item 该分量定义为$d_n^f$
    \item 证明$d_0^f$等于f诱导的A的accosiated graded上面的map

\end{enumerate}

\paragraph{}之后的axiom全部删除，按照 Section \ref{ess} 再生成一遍

\begin{thebibliography}{99}


\end{thebibliography}

\end{sloppy}
\end{document}
